% !TEX root = template.tex

\section{Concluding Remarks}
\label{sec:conclusions}

In this project we tried to tackle the problem of music generation by expanding the ideas already presented in \cite{Yang-2017}. We introduced a new way of setting the adversarial game, by using Wassersetin distance and gradient penalty instead of BCE loss. Furthermore, we gave a more profound structure to both the networks, providing them with a better understading of the music generation process.
We think that this network could be trained on different datasets and still be able to produce meaningful results, due to the various architectural constraints we imposed to the geneartion process.
However, we also think that longer trainings would be needed to fully understand how the network behaves and how different parameter tunings could lead to better results.

What we realised by designing this network is that it is difficult to understand how the different blocks interact with each other and hence how to improve the model.
However, we understood that cloesly analyzing the losses and the results are the most important factors in undesrtanding which are the required changes that should be made to the model.
Each block we added had its own purpose and we could only evaluate its effectiveness by looking at these data.